

\section{La question de la Prononciation : Vivante contre érasmienne}

Dans le vaste chantier de la refondation cognitive de l'enseignement du grec ancien, la question de
la prononciation, loin d'être un point de détail philologique ou une coquetterie d'esthète,
constitue la clé de voûte de l'acquisition langagière. La section III.4.\cite{ref1} de cette thèse
se propose d'examiner une anomalie historique et pédagogique persistante : la domination quasi
exclusive, dans le monde académique occidental, d'une prononciation dite \og érasmienne\fg{}, dont
le caractère artificiel et désincarné entrave les mécanismes neuropsychologiques fondamentaux de
l'apprentissage.

L'enseignement traditionnel du grec ancien, tel qu'il s'est cristallisé au XIXe siècle en France et
en Europe, repose sur un paradoxe. Alors que les sciences cognitives et la didactique des langues
vivantes (SLA \- \textit{Second Language Acquisition}) ont démontré le rôle central de la boucle
phonologique et de l'immersion auditive dans la mémorisation et la fluidité de lecture, le grec est
enseigné comme un code visuel muet, ou pire, oralisé selon des conventions fluctuantes qui varient
d'une nation à l'autre, voire d'une salle de classe à l'autre. Cette \og cacophonie\fg{} académique,
héritée d'un malentendu humaniste et figée par l'institution scolaire, prive le cerveau de
l'apprenant des ancrages sensoriels nécessaires à l'intériorisation de la langue.

Ce rapport de recherche se livre à une analyse exhaustive des dynamiques historiques, linguistiques
et cognitives en jeu. Il s'agira de déconstruire le mythe érasmien, d'en exposer les failles
pédagogiques à la lumière des neurosciences de l'éducation, et d'évaluer la pertinence des
alternatives \og vivantes\fg{} qui émergent aujourd'hui : la prononciation historique restituée
(Koine impériale) et la prononciation moderne (Reuchlinienne). L'enjeu est de taille : il s'agit de
passer d'une philologie du déchiffrement à une pédagogie de la parole, seule capable de \og
ressusciter\fg{} cognitivement le grec ancien.



\subsection{Archéologie d'un dogme : genèse, réception et dérive de la prononciation érasmienne}

Pour comprendre la résistance opiniâtre du modèle érasmien dans les institutions universitaires, il
est impératif d'en retracer la généalogie. Ce système ne s'est pas imposé par la seule force de sa
vérité scientifique, mais à travers une série de controverses culturelles, politiques et
idéologiques qui ont progressivement coupé le grec de ses racines vivantes.



\subsection{Le \textit{dialogus} de 1528 : entre érudition humaniste et canular savant?}

La pierre angulaire de la prononciation scolaire actuelle est le traité de Desiderius Erasmus,
\textit{De recta Latini Graecique sermonis pronuntiatione Dialogus}, publié à Bâle en 1528 chez
Froben.\cite{ref1} Il est crucial de noter la nature même de ce texte : loin d'être un manuel aride
de phonétique, il s'agit d'un dialogue plein d'esprit (\textit{witty Erasmian style}), mettant en
scène des animaux — un lion et un ours — pour discuter, et souvent moquer, les \og barbarismes\fg{}
des prononciations vernaculaires de l'époque.\cite{ref1}

L'intention d'Érasme était complexe. D'une part, il s'agissait d'une véritable tentative de
reconstruction philologique, s'appuyant sur les autorités antiques comme Quintilien pour retrouver
la distinction des quantités vocaliques (longues/brèves) et la nature musicale de l'accent, perdues
dans la prononciation byzantine unifiée.\cite{ref1} D'autre part, le projet participait d'une utopie
humaniste : créer une \textit{koinè} savante paneuropéenne, un dialecte international standardisé
pour la \og République des Lettres\fg{}, libéré des particularismes nationaux qui fragmentaient le
latin et le grec.\cite{ref1}

Cependant, une ombre plane sur la genèse de ce texte. Des chercheurs contemporains, reprenant des
critiques anciennes (comme celles de Vossius en 1635), soulèvent l'hypothèse que le
\textit{Dialogus} ait pu trouver son origine dans un canular. Chrys Caragounis, dans son analyse
critique \textit{The Error of Erasmus}, rapporte que l'humaniste aurait été piégé par un visiteur
suisse, Henricus Glareanus, qui lui aurait raconté avoir entendu des Grecs prononcer leur langue
avec des sons distincts pour chaque voyelle (notamment le \textit{êta} comme un /e/ et non un /i/).
Érasme, séduit par cette idée qui flattait son intellect, aurait rédigé son dialogue pour
s'attribuer la paternité de cette \og découverte\fg{}, avant de réaliser la supercherie mais sans
pouvoir se dédire.\cite{ref4} Que cette anecdote soit factuelle ou apocryphe, elle souligne un point
fondamental : le système érasmien est né d'une construction intellectuelle occidentale, en rupture
consciente avec la tradition orale ininterrompue des locuteurs natifs grecs.



\subsection{La guerre des prononciations : érasmiens contre reuchliniens}

La publication du \textit{Dialogus} a ouvert une \og vaste querelle\fg{} qui a traversé les siècles,
opposant deux camps irréconciliables : les \og Érasmiens\fg{}, partisans d'une reconstruction
savante, et les \og Reuchliniens\fg{} (du nom de Johannes Reuchlin), défenseurs de la prononciation
traditionnelle byzantine/moderne, apprise auprès des exilés grecs après la chute de
Constantinople.\cite{ref6}

Ce conflit n'était pas seulement linguistique ; il était éminemment politique. En France, au XIXe
siècle, la question de la prononciation du grec s'est trouvée intimement liée au mouvement du
\textbf{Philhellénisme} et à la Guerre d'Indépendance grecque (1821-1829).\cite{ref3}



\subsection{Le "schisme philologique" en france}

Sous le Premier Empire et la Restauration, des figures comme \textbf{Fleury de Lécluse}, premier
titulaire de la chaire de littérature grecque à la Faculté des Lettres de Toulouse, se sont érigées
contre l'érasmisme. Dans sa \textit{Dissertation sur la prononciation grecque} (1829), Lécluse
argumente avec véhémence en faveur de la prononciation \og orientale\fg{} ou nationale.\cite{ref3}
Ses arguments résonnent avec une modernité frappante pour notre thèse de refondation cognitive :

\begin{itemize}\item \textbf{Souveraineté Linguistique :} Lécluse dénonce l'arrogance des savants occidentaux (Français, Allemands, Hollandais) qui prétendent apprendre aux Grecs comment prononcer leur propre langue. Il compare cette attitude à celle d'un étranger qui viendrait dicter aux Français la \og vraie\fg{} prononciation de Racine ou Molière sur la base de théories abstraites.\cite{ref3}\item \textbf{Vitalité et Communication :} Pour les philhellènes, adopter la prononciation moderne était un acte de solidarité et de pragmatisme. Il s'agissait de rétablir des \og rapports intimes\fg{} avec le peuple grec régénéré, de traiter le grec comme une langue vivante utile aux échanges diplomatiques et culturels, et non comme un cadavre disséqué.\cite{ref3}\item \textbf{Preuves Historiques :} Lécluse invoque Rabelais (Pantagruel, Livre II, ch. 9) pour prouver qu'à la Renaissance encore, la prononciation \og moderne\fg{} (iotacisante) était la norme acceptée avant l'imposition de la réforme érasmienne. Il montre que Rabelais transcrivait le grec en caractères latins selon la phonétique moderne (ex: \textit{emous} pour \textit{emou}, \textit{panta} pour \textit{panta}).\cite{ref8}\end{itemize}

\subsection{Le désenchantement et la victoire de l'académie}

Malgré la vigueur de ces arguments, le courant érasmien l'a emporté, porté par
l'institutionnalisation de l'enseignement secondaire et universitaire sous la IIIe République. Le
\og désenchantement\fg{} politique vis-à-vis de la Grèce moderne, perçue par certains élites comme
une nation \og dégénérée\fg{} indigne de ses ancêtres antiques, a favorisé une \og captation
d'héritage\fg{}.\cite{ref3} L'Université française, sous l'influence de réformateurs comme
\textbf{Michel Bréal} et dans le sillage des instructions officielles de 1880 initiées par Jules
Ferry, a scellé le sort du grec.\cite{ref3}

La réforme de 1880, en cherchant à rationaliser et \og scientifiser\fg{} l'enseignement, a entériné
la rupture. Le grec est devenu un \og monument\fg{}, un objet d'étude fermé sur lui-même, dont la
prononciation devait servir à l'analyse grammaticale et orthographique (distinguer \textit{êta} de
\textit{iota} pour l'orthographe) plutôt qu'à la communication. Cette \og monumentalisation\fg{} a
eu pour effet pervers de transformer la langue en un exercice de décodage visuel, évacuant la
dimension sonore et vivante qui est pourtant le moteur de l'apprentissage naturel.



\subsection{De la "restitution" à l'artifice : anatomie de la babel érasmienne}

Le terme \og érasmien\fg{} suggère une unité de méthode et une rigueur scientifique. Or, l'analyse
de la pratique réelle dans les universités occidentales révèle une réalité tout autre : une
cacophonie de systèmes nationaux contradictoires, que Chrys Caragounis qualifie d'erreur historique
majeure.



\subsection{L'illusion de l'universalité : les "érasmismes" nationaux}

Loin de l'idéal d'Érasme d'une prononciation paneuropéenne, l'érasmisme s'est fragmenté en dialectes
académiques locaux, influencés par la phonologie de la langue maternelle des apprenants. Il n'y a
pas \textit{une} prononciation érasmienne, mais des \og érasmismes\fg{} anglais, français, allemands
ou américains, souvent mutuellement inintelligibles.\cite{ref11}

Le tableau ci-dessous illustre ces divergences pour quelques graphèmes clés, démontrant
l'incohérence du système \og standard\fg{} actuel :

\begin{table}[h!]\small\centering\begin{tabular}{| l | l | l | l | l |}\hline\textbf{Graphème} & \textbf{Grec Ancien (Attique Ve s. Restitué)} & \textbf{Grec Moderne (Reuchlinien)} & \textbf{Érasmien \og Français\fg{} (Usage scolaire)} & \textbf{Érasmien \og Anglais/US\fg{} (Usage scolaire)} \\ \hline\textbf{β (Bêta)} & /b/ (occlusive) & /v/ (fricative) & /b/ & /b/ \\ \hline\textbf{γ (Gamma)} & /g/ (occlusive) & /ɣ/ ou /j/ (fricative) & /g/ (dur) & /g/ (dur) \\ \hline\textbf{δ (Delta)} & /d/ (occlusive) & /ð/ (fricative \og th\fg{}) & /d/ & /d/ \\ \hline\textbf{θ (Thêta)} & /tʰ/ (aspirée) & /θ/ (fricative \og th\fg{}) & /t/ & /θ/ (\og thin\fg{}) \\ \hline\textbf{φ (Phi)} & /pʰ/ (aspirée) & /f/ (fricative) & /f/ & /f/ \\ \hline\textbf{χ (Chi)} & /kʰ/ (aspirée) & /x/ ou /ç/ (fricative) & /k/ & /k/ ou /kʰ/ \\ \hline\textbf{η (Êta)} & /ɛː/ (long ouvert) & /i/ & /ɛ/ (comme \og fête\fg{}) & /eɪ/ (comme \og late\fg{}) \\ \hline\textbf{υ (Upsilon)} & /y/ (u français) & /i/ & /y/ (u français) & /u/ (\og boot\fg{}) ou /ju/ \\ \hline\textbf{ει (Diphtongue)} & /eː/ (puis /iː/) & /i/ & /ɛj/ (\og soleil\fg{}) & /aɪ/ (\og eye\fg{}) ou /eɪ/ \\ \hline\textbf{οι (Diphtongue)} & /oi/ & /i/ & /wa/ (\og loi\fg{}) & /ɔɪ/ (\og boy\fg{}) \\ \hline\textbf{ου (Diphtongue)} & /uː/ & /u/ & /u/ (\og ou\fg{}) & /aʊ/ (\og out\fg{}) ou /u/ \\ \hline\end{tabular}\end{table}(Sources des variations phonétiques : 11)

\begin{itemize}\item \textbf{Le Cas Français :} L'érasmisme français applique systématiquement les règles phonétiques du français moderne. Les accents toniques grecs sont ignorés au profit d'une oxytonisation (accentuation sur la dernière syllabe) typique du français. La diphtongue \textbf{οι} est prononcée /wa/ (comme dans \og loi\fg{}), une aberration historique totale qui n'a jamais existé en grec antique (où c'était /oi/ puis /y/ puis /i/). De même, \textbf{au} et \textbf{eu} sont souvent monophtongués.\cite{ref16}\item \textbf{Le Cas Anglo-Saxon :} Les anglophones introduisent des glides et des diphtongues parasites (le \textit{êta} devient /eɪ/), et appliquent un accent d'intensité très fort qui détruit la structure quantitative du vers grec. Le \textbf{χ} (Chi) est souvent prononcé comme un simple /k/, perdant toute distinction avec le Kappa.\cite{ref13}\end{itemize}Cette \og Babel\fg{} a des conséquences désastreuses : lors de colloques internationaux, les
spécialistes du Nouveau Testament ou de la littérature classique sont souvent incapables de se
comprendre lorsqu'ils lisent une citation en grec, devant recourir à la traduction écrite ou à
l'anglais pour communiquer.\cite{ref13} L'érasmisme a échoué dans sa mission première de
communication savante.



\subsection{L'artifice comme obstacle cognitif}

Plus grave encore que l'inexactitude historique est l'artificialité cognitive du système.
L'érasmisme scolaire traite le grec comme une algèbre. Il force l'apprenant à produire des sons qui
ne correspondent à aucune réalité linguistique cohérente.

\begin{itemize}\item \textbf{La primauté de l'œil sur l'oreille :} L'érasmisme justifie ses choix (comme prononcer \textbf{ει} différemment de \textbf{ι}) par la nécessité de distinguer l'orthographe. C'est une pédagogie de l'écrit qui nie la réalité de l'évolution de la langue, où l'homophonie est naturelle (comme en français entre \og verre\fg{}, \og vert\fg{}, \og vers\fg{}).\item \textbf{L'instabilité phonologique :} Faute de modèles natifs ou d'enregistrements standards, l'étudiant \og bricole\fg{} sa prononciation. Cette instabilité empêche la fixation des traces mnésiques. Un mot dont la forme sonore est floue est un mot qui ne s'ancre pas dans la mémoire à long terme.\end{itemize}Comme le souligne Chrys Caragounis, l'érasmisme est une \og prononciation artificielle\fg{} qui va à
l'encontre de la nature même du langage, transformant l'apprentissage en une tâche de décryptage
laborieuse plutôt qu'en une acquisition de compétences.\cite{ref4}



\subsection{Les fondements cognitifs de l'oralité : pourquoi le cerveau a besoin de son}

La thèse de la \og refondation cognitive\fg{} postule que les méthodes d'enseignement doivent
s'aligner sur le fonctionnement réel du cerveau. Or, les recherches en acquisition des langues
secondes (SLA \- \textit{Second Language Acquisition}) et en neuropsychologie condamnent sans appel
l'approche muette ou artificiellement oralisée du grec ancien.



\subsection{La boucle phonologique et la mémoire de travail}

Le modèle de la mémoire de travail de Baddeley \& Hitch identifie la \textbf{boucle phonologique}
(\textit{phonological loop}) comme un composant essentiel de l'apprentissage linguistique. Cette
boucle permet de maintenir temporairement l'information verbale par un processus de répétition
subvocale (la \og petite voix\fg{} dans la tête).\cite{ref19}

Les recherches démontrent que :

1. \textbf{L'acquisition du vocabulaire dépend de la boucle phonologique :} La capacité à apprendre
de nouveaux mots (formes sonores) est directement corrélée à la capacité de répétition
phonologique.\cite{ref21} Si la prononciation est incertaine, incohérente ou purement visuelle
(comme c'est souvent le cas avec l'érasmienne scolaire), la boucle phonologique ne peut pas
fonctionner efficacement. L'encodage en mémoire à long terme échoue.

2. \textbf{Séquençage et Rétention :} L'apprentissage de la grammaire (morphologie) nécessite la
détection de régularités sonores (statistiques). Un input auditif riche et constant permet au
cerveau d'extraire ces règles implicitement. L'approche érasmienne, pauvre en input auditif (limitée
à la lecture de phrases isolées), prive le cerveau de ce mécanisme d'apprentissage
statistique.\cite{ref21}



\subsection{Subvocalisation et fluence de lecture}

Contrairement à une idée reçue, la lecture silencieuse n'est pas purement visuelle. Elle implique
une \textbf{subvocalisation} active : le cerveau \og prononce\fg{} les mots intérieurement pour en
extraire le sens.\cite{ref22}

\begin{itemize}\item \textbf{Le Coût de l'Artifice :} Un apprenant formé à l'érasmienne, qui doit déchiffrer lettre par lettre sans flux prosodique naturel, subit une surcharge cognitive. Sa subvocalisation est lente et hachée. Cela sature la mémoire de travail : avant d'arriver à la fin d'une phrase complexe de Platon ou de Saint Paul, il a déjà oublié le début.\cite{ref24}\item \textbf{L'Avantage de la Prononciation Vivante :} Une prononciation standardisée et fluide (qu'elle soit moderne ou restituée de manière immersive) permet une subvocalisation rapide et automatique. Cela libère les ressources attentionnelles pour les processus de haut niveau : analyse syntaxique, compréhension sémantique et appréciation esthétique.\cite{ref25}\end{itemize}

\subsection{L'importance de l'input auditif et de l'immersion}

Les théories modernes du SLA (Krashen, input hypothesis) et les neurosciences confirment que le
cerveau est câblé pour apprendre le langage par l'oreille avant l'œil. L'absence de stimuli auditifs
dans l'enseignement classique du grec (réduit à la traduction) explique les taux d'échec élevés et
l'incapacité de la plupart des étudiants à lire couramment après des années d'étude.\cite{ref26}

Les méthodes qui intègrent une prononciation vivante et beaucoup d'audio (comme celles du Biblical
Language Center ou de l'Institut Polis) montrent des résultats supérieurs en termes de rétention et
de compétence de lecture, car elles activent pleinement les circuits neuronaux du
langage.\cite{ref28}



\subsection{Les alternatives vivantes : modèles historiques et pédagogiques}

Face à l'échec cognitif de l'érasmisme, la refondation de l'enseignement doit se tourner vers des
systèmes phonologiques cohérents, capables de porter une pédagogie de l'oralité. Deux voies
principales s'offrent aujourd'hui, chacune avec ses justifications historiques et pratiques : le
modèle Reuchlinien (Moderne) et le modèle Restitué (Koine ou Attique).



\subsection{L'option "néo-hellénique" (reuchlinienne) : la continuité vivante}

Cette option consiste à adopter la prononciation du grec moderne standard pour lire tous les textes
anciens. C'est la pratique ininterrompue de l'Église Orthodoxe et du système éducatif
grec.\cite{ref6}

\begin{itemize}\item \textbf{Richesse des Ressources Auditives :} C'est le seul système qui offre un accès immédiat à une communauté vivante de locuteurs et à des milliers d'heures d'enregistrements (liturgie, Nouveau Testament lu par des natifs).\cite{ref29} Pour la boucle phonologique, c'est un atout inestimable.\item \textbf{Simplicité Cognitive :} Le système vocalique moderne est simple (5 voyelles /a, e, i, o, u/), ce qui réduit la charge cognitive initiale pour l'apprenant.\cite{ref6}\item \textbf{Unité Historique :} Elle souligne la continuité de la langue grecque sur 3000 ans, combattant l'idée psychologiquement paralysante d'une \og langue morte\fg{}.\cite{ref33}\end{itemize}\begin{itemize}\item \textbf{Le problème de l'Iotacisme :} La fusion de six graphèmes (ι, η, υ, ει, οι, υι) en un seul son /i/ est l'argument principal des érasmiens, qui craignent la confusion orthographique (ex: \textit{humeis} \og vous\fg{} vs \textit{hēmeis} \og nous\fg{}, prononcés identiquement /imis/).\cite{ref13}\item \textit{Réponse :} Les recherches montrent que les homophones contextuels ne gênent pas la compréhension globale (comme en anglais ou français). De plus, l'apprentissage par immersion permet de distinguer les mots par le contexte syntaxique.\cite{ref34}\end{itemize}

\subsection{L'option "koine restituée" (modèle buth/rico/ranieri) : la voie médiane}

Une \og troisième voie\fg{} gagne du terrain dans les cercles de pédagogie active (\textit{Living
Greek movement}), portée par des linguistes comme Randall Buth (\textit{Biblical Language Center}),
Christophe Rico (\textit{Polis Institute}) et Luke Ranieri.\cite{ref28} Il s'agit de reconstruire la
prononciation de la \textbf{Koine Impériale} (Ier s. av. J.-C. – IVe s. ap. J.-C.).

Principes Phonologiques :

Ce système se distingue à la fois de l'Attique classique et du Moderne :

\begin{itemize}\item \textbf{Fricativisation des Consonnes :} Comme en moderne, φ, θ, χ sont prononcés comme des fricatives (/f/, /θ/, /x/), ce qui est plus naturel pour les apprenants européens que les aspirées antiques.\cite{ref14}\item \textbf{Conservation du /y/ (Upsilon/Oi) :} C'est la différence cruciale avec le moderne. Dans la Koine de cette période, \textbf{υ} et \textbf{οι} se prononçaient encore /y/ (comme le 'u' français ou 'ü' allemand) et ne s'étaient pas encore confondus avec /i/.\cite{ref14}\item \textbf{Distinction Êta/Iota :} Selon les variantes (Lucian de Ranieri vs Koine de Buth), le \textit{êta} peut être conservé comme un /e/ fermé long ou déjà fusionné vers /i/.\cite{ref37}\end{itemize}1. \textbf{Désambiguïsation Pédagogique :} En gardant le son /y/, ce système résout les homophonies
les plus gênantes (notamment \textit{humeis} / \textit{hēmeis}), facilitant l'acquisition de
l'orthographe sans sacrifier la cohérence historique.\cite{ref14}

2. \textbf{Pont Historique :} C'est la prononciation réelle des auteurs du Nouveau Testament, de
Plutarque, de Galien et des premiers Pères. Elle permet une lecture authentique de ce corpus
immense.

3. \textbf{Support Méthodologique :} Ce modèle est soutenu par des méthodes communicatives modernes
(audio, vidéo, TPR) qui garantissent un apprentissage vivant.\cite{ref27}



\subsection{L'option "attique restituée" (allen/stratakis) : l'idéal scientifique}

Défendue par W. Sidney Allen (Vox Graeca) et mise en pratique par des récants comme Ioannis
Stratakis, cette prononciation vise l'exactitude du Ve siècle av. J.-C. (tonalité, distinction
longues/brèves, aspirées).\cite{ref6}

Bien que scientifiquement \og pure\fg{} pour lire Platon ou Sophocle, elle impose une charge
cognitive extrêmement lourde (maîtrise de l'accent tonal musical). Pour une initiation ou une
refondation cognitive visant la fluidité, elle risque d'être un obstacle, sauf pour des étudiants
avancés en poésie.\cite{ref42}



\subsection{Synthèse comparative et recommandations pour la thèse}

Afin d'éclairer le choix pédagogique dans le cadre de la refondation cognitive, le tableau suivant
synthétise les différences phonologiques critiques et leur impact sur l'apprentissage.



\subsection{Tableau comparatif des systèmes phonologiques}

\begin{table}[h!]\small\centering\begin{tabular}{| l | l | l | l | l |}\hline\textbf{Caractéristique} & \textbf{Grec Ancien (Attique Ve s. Restitué)} & \textbf{Koine Restituée (Impériale)} & \textbf{Grec Moderne (Reuchlinien)} & \textbf{Érasmien \og Scolaire\fg{} (Variable selon pays)} \\ \hline\textbf{Objectif} & Exactitude historique (Ve s. av. JC) & Pédagogie active \& Historicité (Ier-IVe s.) & Continuité culturelle \& Ressources vivantes & Convention orthographique (lecture visuelle) \\ \hline\textbf{Voyelle η (Êta)} & /ɛː/ (long ouvert \og bête\fg{}) & /e/ (fermé) ou /i/ & \textbf{/i/} & /ɛ/ (fr) ou /eɪ/ (anglais) \\ \hline\textbf{Voyelle υ (Upsilon)} & \textbf{/y/} (u français) & \textbf{/y/} (u français) & \textbf{/i/} & /y/ (fr) ou /u/ (anglais) \\ \hline\textbf{Diphtongue ει} & /eː/ puis /iː/ & \textbf{/i/} & \textbf{/i/} & /ɛj/ ou /aɪ/ \\ \hline\textbf{Diphtongue οι} & /oi/ & \textbf{/y/} (u français) & \textbf{/i/} & /wa/ (fr) ou /ɔɪ/ (anglais) \\ \hline\textbf{Diphtongue αι} & /ai/ & \textbf{/e/} & \textbf{/e/} & /aj/ ou /aɪ/ \\ \hline\textbf{Consonnes φ, θ, χ} & Aspirées /pʰ, tʰ, kʰ/ & Fricatives /f, θ, x/ & Fricatives /f, θ, x/ & Fricatives ou Occlusives (incohérent) \\ \hline\textbf{Accentuation} & Tonal (Hauteur musicale) & Stress (Intensité) & Stress (Intensité) & Stress (souvent mal placé ou ignoré) \\ \hline\textbf{Charge Cognitive} & Très Haute (difficile) & Moyenne (équilibrée) & Basse (simple) & Haute (car incohérente/artificielle) \\ \hline\end{tabular}\end{table}(Sources compilées : 13)



\subsection{Vers une pragmatique de la prononciation}

Dans le cadre de la thèse, la recommandation ne doit pas être dogmatique mais pragmatique, orientée
vers l'efficacité cognitive (l'aisance de l'apprenant).

1. \textbf{Rejet de l'Érasmisme Artificiel :} Il est impératif d'abandonner la prononciation
érasmienne traditionnelle, surtout dans ses variantes nationales (comme le /wa/ français pour
\textbf{οι}). Elle constitue un bruit cognitif qui empêche l'immersion et isole l'étudiant
francophone de la communauté internationale et historique.\cite{ref4}

2. \textbf{Adoption de la Koine Restituée pour l'Enseignement Général :} Le modèle de la
\textbf{Koine Impériale (type Buth/Rico)} apparaît comme le compromis optimal pour une refondation
cognitive.\cite{ref28}

\begin{itemize}\item Elle permet l'usage de méthodes communicatives (parler en classe).\item Elle offre une distinction phonologique suffisante (/y/ vs /i/) pour soutenir l'orthographe sans la complexité de l'Attique archaïque.\item Elle est historiquement légitime pour une grande partie du corpus chrétien et impérial.\end{itemize}3. \textbf{Valorisation du Moderne pour la Continuité :} Pour les parcours axés sur la théologie
orthodoxe ou l'histoire byzantine, le modèle Reuchlinien est parfaitement viable et doit être
déstigmatisé. Il offre l'accès le plus direct à la langue \og vivante\fg{}.\cite{ref34}



\subsection{Conclusion}

La \og Question de la Prononciation\fg{} dépasse largement le cadre de la phonétique corrective.
Elle touche au cœur de la relation que l'apprenant entretient avec la langue. Maintenir l'artifice
érasmien, c'est condamner le grec à rester une langue de papier, déchiffrée par l'œil mais muette
pour l'esprit.

La refondation cognitive de l'enseignement du grec ancien exige de redonner corps à la langue. En
adoptant une prononciation vivante — qu'il s'agisse de la Koine Restituée pour sa rigueur
pédagogique ou du Grec Moderne pour sa vitalité culturelle — nous permettons au cerveau de réactiver
ses mécanismes naturels d'acquisition (boucle phonologique, apprentissage statistique, immersion).
C'est à ce prix que le grec ancien pourra cesser d'être un \og monument\fg{} froid pour redevenir
une parole vivante, capable de résonner dans l'intelligence contemporaine.

